% ============================================
% DOC-20
% Extraído de main.pandoc.tex
% ============================================
\fragtitle{doc}{D}{DOC}{20}{Termo de Cadeia de Custódia}
\markboth{DOC-20 --- Termo de Cadeia de Custódia}{}

\noindent \textit{Documento extraído do arquivo pericial.}

\vspace{0.5em}\noindent\rule{\textwidth}{0.5pt}\vspace{0.5em}

\noindent {\footnotesize\texttt{ARQUIVO P\'UBLICO $\cdot$ PROVA DOCUMENTAL $\cdot$ PROCESSO 1.260}}\\[-0.2em]

\noindent \textbf{Instituto Médico-Legal de Barbacena --- Documento de Controle}\\
\noindent \textbf{Objeto:} caderno encadernado em couro, dimensões 16\,cm × 22\,cm, aproximadamente 200 páginas, sem numeração original. \\
\noindent \textbf{Atribuído a:} Joaquim Antônio Correia (Paciente 1.260). \\
\noindent \textbf{Natureza:} diário pessoal / registro clínico não oficial.

\underline{Registro de movimentação do objeto:}

\begin{moltablefive}
\hline
\noindent \textbf{N\textsuperscript{o}} & \textbf{Data} & \textbf{Origem} & \textbf{Destino} & \textbf{Responsável} \\ \hline
\noindent 1 & 14/jun/1979 & Cova 1.260, Cemitério do Colônia & IML --- Necropsia & Genésio Barbosa (coveiro) \\
\noindent 2 & 14/jun/1979 & IML --- Necropsia & IML --- Cofre & Dr. Antônio Tavares (legista) \\
\noindent 3 & 16/jun/1979 & IML --- Cofre & \textit{[s/ preenchimento]} & Dr. Heitor Oliveira \\
\noindent 4 & \textit{[ilegível]} & \textit{[não registrado]} & Arq. Público Mineiro & \textit{[s/ identificação]} \\
\noindent 5 & 12/jan/2026 & APM --- Consulta & APM --- Pesquisa & Dra. Lúcia Mendes \\
\noindent 6 & 28/jul/2026 & APM --- Pesquisa & APM --- Cofre & Dra. Lúcia Mendes \\
\noindent 7 & \textit{[em branco]} & \textit{[em branco]} & \textit{[em branco]} & \textit{[em branco]} \\
\hline
\end{moltablefive}

\noindent \textbf{Observações anexadas:}

\noindent \textit{1. A entrada nº 3 (Dr. Heitor Oliveira) não possui autorização formal de retirada. O Livro de Registro do IML registra o nome do médico, mas não há assinatura de testemunha, carimbo de autoridade competente ou justificativa por escrito.}
\noindent \textit{2. A entrada nº 4 foi reconstituída a partir de menção indireta em correspondência interna do Arquivo Público Mineiro. A data original e o nome do responsável pela transferência foram removidos do documento físico, deixando apenas um retângulo vazio onde antes havia texto. A análise óptica sugere que o nome foi raspado com lâmina.}
\noindent 3. O campo nº 7 foi deixado em branco. Em 2026, quando o objeto foi transferido ao cofre definitivo, o funcionário responsável pelo preenchimento — cujo nome não consta no documento — declarou verbalmente que ``não cabia a ele preencher a última linha''. Indagado sobre o motivo, respondeu: \textquotedblleft{}\textit{Porque o diário não vai ficar aqui.}\textquotedblright{}

\noindent 4. Ao ser questionado se sabia para onde o objeto seria levado em seguida, o funcionário hesitou e disse: \textquotedblleft{}\textit{Para a mão de quem está esperando.}\textquotedblright{}

\noindent 5. Em declaração anexa, o diretor Ronaldo Torres informou que guardou o objeto entre 2019 e 2026 e que o devolveu à recepção em 13/05/2026, sem se identificar.

\noindent 6. A exumação de agosto de 1979, que recuperou o caderno, não consta desta tabela; o registro foi reconstituído a partir do depoimento do coveiro (DOC-01).

\noindent \textbf{Conclusão do arquivista:}
\noindent A cadeia de custódia do objeto não é contínua. Há duas lacunas documentais não preenchidas e uma entrada cujo responsável não pode ser confirmado. O objeto circulou por pelo menos 47 anos sem registro formal de paradeiro.
\noindent Este termo é, ele mesmo, um documento de ausências.

\noindent \textit{Nota do Arquivista: os fragmentos desta parte foram organizados por Dra. Lúcia Mendes durante sua pesquisa em 2026. As anotações dela estão em itálico nas margens dos documentos originais. O documento mais antigo do arquivo --- de 1853 --- foi encontrado por ela e anexado ao processo.}

\vspace{0.5em}
\noindent\textit{\textbf{↪ Rotas:} \rota{DOC-01} \rotasep \rota{LUC-02}}